\section{Extending prescribed maps through laminar flags}
A family of subsets is laminar if any two are nested or disjoint.
The following statement retains an arbitrary prescribed local map;
that is the part not stated in the original UB.F interface.

\begin{theorem}[UB.Laminar]\label{thm:laminar}
Let $\Omega$ be finite, let $Q\subseteq\Omega$, and let
$(A_i)_{i\in I}$ be a finite indexed laminar family of subsets of
$\Omega\setminus Q$. Choose $J\subseteq I$ and subsets $X_i$,
$i\in J$, with
\[
 |X_i|=|A_i|,\qquad |X_i\cap X_j|=|A_i\cap A_j|.
\]
Write $D=\Omega\setminus\bigcup_{i\in J}X_i$.
Every injection $u:Q\to D$ extends to a permutation $g$ of $\Omega$
satisfying $gA_i=X_i$ for all $i\in J$.
\end{theorem}
\begin{proof}
If a selected set is $\Omega$, then $Q=\varnothing$ and its target
is necessarily $\Omega$. Discard that vacuous constraint; the prescribed
map is empty in this case. We may therefore assume that every selected
set is proper, using the complement of the remaining targets as the
root region.
The intersection cardinalities enforce the same disjointness and
containment relations among the $X_i$ as among the $A_i$. Equal
standard sets give equal target sets and can be treated once. Adjoin
$\Omega$ above all the selected sets. For each selected container,
subtract its maximal selected proper subcontainers. Since those
subcontainers are disjoint, the difference has size equal to the
container's size minus the sum of their sizes. The corresponding
standard and target differences therefore have equal cardinalities.
They partition $\Omega$ on each side. The root difference contains
$Q$ on the standard side and is exactly $D$ on the target side.
Extend $u$ to a bijection of these two differences, and choose arbitrary
bijections on every other corresponding pair of differences. Their
union is the required permutation.
\end{proof}


\begin{lemma}[UB.Character]\label{lem:character}
Suppose at least one permutation realizes prescribed laminar images and
a prescribed partial bijection. Fix one such permutation $g_0$, and
let $K$ be the subgroup preserving every standard selected subset and
fixing every point in the domain of that partial bijection.
All realizations are exactly $g_0K$. For a homomorphism
$\chi:S_\Omega\to C$ and $c\in C$, a realization with $\chi(g)=c$
exists exactly when $\chi(g_0)^{-1}c\in\chi(K)$.
\end{lemma}
\begin{proof}
Two realizations differ on the right by a permutation satisfying exactly
the defining conditions of $K$, and every such right multiplication
preserves the prescribed data. Applying $\chi$ proves the second claim.
\end{proof}


These statements make explicit the extension freedom in the familiar atom-matching proof. Their use below does not claim a new historical classification or a prior-art determination.
